\documentclass[10pt]{article}
\usepackage[margin=1in]{geometry}
\usepackage{amsmath,amssymb,mathtools}
\usepackage{booktabs}
\usepackage{xcolor}
\usepackage[hidelinks]{hyperref}
\usepackage{microtype}
\usepackage{fancyhdr}
\definecolor{navy}{RGB}{24,55,95}
\pagestyle{fancy}
\setlength{\headheight}{14pt}
\fancyhf{}
\lhead{Alternating groups}
\rhead{Exact finite-$n$ proof and verification}
\cfoot{\thepage}
\newcommand{\Sym}{\operatorname{Sym}}
\newcommand{\M}{\mathcal M}
\title{\color{navy}\textbf{The Exact Finite-$n$ $\sigma$-MGF for $A_n$}\\[2mm]
\large Proof, orbit splitting, and matrix verification}
\author{Computational proof companion}
\date{17 July 2026}

\begin{document}
\maketitle
\begin{abstract}
For $n\ge2$, this note proves the exact cycle-type formula for the permutation
representation of $A_n$, derives the correction to the $S_n$ vector-partition
count from orbit splitting, and proves equality with $S_n$ in degree at most
$n-2$.  Explicit even permutation matrices and Reynolds projectors verify the
formula, including a boundary case with a nonzero splitting correction.
\end{abstract}

\section{From \texorpdfstring{$S_n$}{S-n} to \texorpdfstring{$A_n$}{A-n}}
For any class function $F$ on $S_n$, the indicator of an even permutation gives
\[
 \mathbb E_{A_n}[F]=\frac1{n!}\sum_{g\in S_n}(1+\operatorname{sgn}(g))F(g).
\]
As in the symmetric-group calculation, the coefficient of $m_\tau$ is the
trace, hence the rank, of the Reynolds projector on
$W_\tau=\bigotimes_b\Sym^{\tau_b}(\mathbb C^n)$.

If $\lambda\vdash n$, then
\[
 \operatorname{sgn}(\lambda)=(-1)^{n-\ell(\lambda)},\qquad
 z_\lambda=\prod_{d\ge1}d^{m_d(\lambda)}m_d(\lambda)!.
\]
Combining the sign indicator with the determinant of each permutation cycle
proves the exact expression
\[
 \boxed{\M_{A_n}=\sum_{\lambda\vdash n}
 \frac{1+(-1)^{n-\ell(\lambda)}}{z_\lambda}
 \prod_i\prod_{d\ge1}(1-t_i^d)^{-m_d(\lambda)}.}
\]
Only even cycle types remain, each with coefficient $2/z_\lambda$.

\section{Generating function}
Set
\[
 G_+(u)=\prod_{\mathbf v\ge0}(1-ut^{\mathbf v})^{-1},\qquad
 G_-(u)=\prod_{\mathbf v\ge0}(1+ut^{\mathbf v}).
\]
The first product records multisets of labels.  The second records sets of
distinct labels and is the sign-weighted cycle-index contribution.  Hence, for
$n\ge2$,
\[
 \boxed{\M_{A_n}=[u^n]\bigl(G_+(u)+G_-(u)\bigr).}
\]
Equivalently, applying the exponential formula to both factors gives
\begin{align*}
\M_{A_n}=[u^n]\Bigg[&\exp\left(\sum_{d\ge1}\frac{u^d}{d}
 \prod_i\frac1{1-t_i^d}\right)\\
&+\exp\left(\sum_{d\ge1}\frac{(-1)^{d-1}u^d}{d}
 \prod_i\frac1{1-t_i^d}\right)\Bigg].
\end{align*}

\section{Exact orbit interpretation}
An $S_n$-orbit of monomial-basis labelings is determined by multiplicities
$(a_{\mathbf v})_{\mathbf v\ge0}$.  Its stabilizer is
$\prod_{\mathbf v}S_{a_{\mathbf v}}$.  If some $a_{\mathbf v}\ge2$, the
stabilizer contains the odd transposition of two equally labeled points, so the
$S_n$-orbit remains a single $A_n$-orbit.  If every $a_{\mathbf v}$ is zero or
one, the stabilizer lies in $A_n$, and the orbit splits into two $A_n$-orbits.
Thus
\begin{align*}
[m_\tau]\M_{A_n}
={}&\#\left\{(a_{\mathbf v})\ge0:\sum a_{\mathbf v}=n,\ 
 \sum a_{\mathbf v}\mathbf v=\tau\right\}\\
&+\#\left\{(a_{\mathbf v})\in\{0,1\}:\sum a_{\mathbf v}=n,\ 
 \sum a_{\mathbf v}\mathbf v=\tau\right\}.
\end{align*}
The first term is the $S_n$ coefficient.  The correction counts sets of
distinct nonzero vectors summing to $\tau$, with either $n$ such vectors or
$n-1$ such vectors and one zero vector.

\section{Stable agreement with \texorpdfstring{$S_n$}{S-n}}
Any labeling of total degree $|\tau|$ has at most $|\tau|$ nonzero labels.  If
$|\tau|\le n-2$, at least two labels are zero.  Their transposition is odd and
lies in the stabilizer, so no $S_n$-orbit can split.  Therefore
\[
 \boxed{[m_\tau]\M_{A_n}=[m_\tau]\M_{S_n}
 \quad\text{whenever }|\tau|\le n-2.}
\]
Outside that range the distinct-label term is the exact correction, not merely
an error estimate.

\section{Verification method and evidence}
The notebook enumerates even permutations, constructs their permutation
matrices, builds the full matrices on each symmetric tensor block, and averages
them to obtain the $A_n$ Reynolds projector.  Its rank is compared independently
with the character average, the even-cycle-type formula, and the $S_n$ orbit
count plus a distinct-label dynamic-programming correction.

\begin{center}
\small
\begin{tabular}{@{}lrrrrrr@{}}
\toprule
Case & $\dim W_\tau$ & rank & cycle & $S_n$ & corr. & total\\
\midrule
$A_3,\tau=(2)$ & 6 & 2 & 2 & 2 & 0 & 2\\
$A_3,\tau=(3)$ & 10 & 4 & 4 & 3 & 1 & 4\\
$A_4,\tau=(2,1)$ & 40 & 4 & 4 & 4 & 0 & 4\\
$A_5,\tau=(2,2)$ & 225 & 9 & 9 & 9 & 0 & 9\\
$A_5,\tau=(3,1)$ & 175 & 7 & 7 & 7 & 0 & 7\\
\bottomrule
\end{tabular}
\end{center}
The $A_3,\tau=(3)$ row deliberately lies outside the stable range and confirms
that one $S_3$ orbit splits, increasing the invariant dimension from three to
four.  The remaining rows test multiblock targets and larger projectors.
Idempotence residuals are at most $1.1\times10^{-15}$.

\paragraph{Conclusion.}
The sign-weighted cycle formula and orbit-splitting formula agree exactly, and
both match direct matrix invariant dimensions in all representative cases.

\end{document}
